\section{AI-assisted methods, verification, and reproducibility}
\label{sec:methods}

This manuscript was developed in a user-directed session with an OpenAI
Codex agent, described by the session as based on GPT-6. The user set
the mathematical objective, requested an editable \LaTeX{} note, and
then asked for substantive new theory integrated with a pedagogical
account of prior results. The model label records the provided session
identity; it is not an independent characterization of the deployed
model or its training.

\paragraph{Tools and organization.}
The agent worked in a local macOS workspace with a zsh shell, local file
editing and Python execution, web retrieval of primary papers, a
\LaTeX{} toolchain, and PDF rendering. It used an author--review--revise
workflow with separate child-agent mandates for structural proofs,
literature comparison, and adversarial examples. A preliminary reviewer
critiqued the research decomposition. The root agent developed the
coverage construction and integrated the manuscript; further reviewers
examined the analytical proofs and actual implementations under
independently framed prompts. The agents shared a filesystem but were
assigned distinct authoring files. The initial research phase involved no publication, submission, or
external correspondence. A subsequent user-directed publication pass
reframed the manuscript around its SAT algorithm and assembled the
manuscript, supporting corpus, and reproducibility workflow as the
public \texttt{understanding-nondeterminism} repository.

A further user-directed expansion developed the manuscript as a deep dive
from computation paths to the counting algorithm. A decomposition critic
reviewed that plan before separate authors drafted the foundations and
worked tables. The root developed the low-excess and conditional-count
consequences, then integrated the sections. Independent mandates reviewed
mathematical accuracy, pedagogical coherence, source attribution, and the
pinning and generation arguments. The revisions made the reading order
follow the worked example and stated exact sampling's probability-one
termination and expected-time contract explicitly.

This organization supplies additional opportunities to detect errors;
it does not turn agreement between agents into mathematical evidence.
The agents can share training-related blind spots. The mathematical
justifications are the complete arguments in the text and the precise
published theorems they invoke, rather than the reviewers' confidence.

\paragraph{How candidate results were developed.}
The initial note combined enumeration, Shannon--Lupanov synthesis,
fixed witness covers, and a forest simulation obtained by fixing shared
vertices. Two limitations guided the next phase. First, a total count
of shared values ignores how they are arranged. This led to the
consistency graph, degree reductions with explicit tables, a cubic
graph, and the pathwidth-based improvement. Second, optimizing witness
coverage separately from sharing charges for computations that can be
reused. This led to selected-restriction occupancies, a correction for
uncovered positive inputs, and a single expected-cost argument.

Each candidate was written with its model and quantifiers before
integration. The literature comparison separated existing elimination,
synthesis, and graph-width tools from the proposed numerical and
distributional consequences. Explicit families then tested what the
new bounds explain that the old guarantee expressions miss. Counterexamples
prevented stronger but false interpretations, such as choosing a
different optimal witness cover independently at each gate.

An earlier focused research phase examined the critical three-way balance.
A finite stability inequality was proved before taking fixed-slack
limits, yielding necessary input, cycle-loss, and width conditions.
An independent construction then realized cubic kernels with nearly
critical supplied budgets. Its AND-only function permitted exact
minimum-size calculations, exposing the difference between the supplied
representation and intrinsic verifier complexity. Separate analytical
review checked the conflict budget, sink construction, equality
splitting, and legal contractions. The final source follow-up retrieved
Nurk's full Russian preprint and additional text of Broering--Lokam;
it narrowed access gaps without establishing priority.

\paragraph{When experiments were useful.}
The user explicitly required that experiments be used only when useful
and never as a substitute for analytical proof of an asymptotic theorem.
The finite checks here serve two purposes: testing circuit-producing
implementations against independent direct evaluation, and checking
exact arithmetic in worked finite instances. There was no empirical
search for a circuit lower bound, no fitted asymptotic exponent, and
no inference of novelty from experiment counts.

The checks found two concrete implementation omissions. The sampled
compiler initially deduplicated full witnesses but not their output
restrictions; this could charge too many output OR gates. It now
deduplicates those restrictions explicitly. A reviewer also found that
the tensor compiler omitted the zero-gate witness-output case already
handled in the proof. That branch and designated-input checks were
added. These findings justify the checks' role without extrapolating
their conclusions beyond their actual domains.

\begin{samepage}
In the exhaustive circuit checks, an encoding consists of one or two
topologically ordered binary truth-table gates, with repeated input
pins allowed and the final gate designated as output. There is one
ordinary and one witness input. This is the precise 9,280-case domain.\par
\end{samepage}

\begin{center}\small
\begin{tabularx}{\linewidth}{@{}>{\raggedright\arraybackslash}p{0.24\linewidth}X@{}}
\toprule
Artifact & Completed finite checks and their scope\\
\midrule
Forest and distribution checks & All 9,280 circuit encodings described
above; 2,000 seeded circuits and three consistency fixtures. Separately,
the exact projection distribution for all 65,536 relations with
$(n,m)=(2,2)$.\\
Coverage compiler & 362 sparse support sets, 896 circuit/sample pairs,
and 140 exact rational occupancy and expected-exception cases.\\
Tensor compiler & Five zero-gate designated-input cases, the 9,280
small circuits, 600 seeded circuits, and 140 symbolic tensor networks
exercising loops, parallel edges, equality splitting, and median ties.\\
Worked families & 25,600 input/witness pairs, support counts for
$k=3,\ldots,10$, finite distribution comparisons, and 141 sparse
correction instances.\\
Kernel realization & 3,729 gate-list and graph constructions,
72,401 legal degree-one/two graph contractions, and exact retained
labeled kernels. Seven representatives also received exhaustive
verifier truth-table and symbolic projection checks.\\
\bottomrule
\end{tabularx}
\end{center}

These enumerations establish agreement of the tested constructions on
their stated finite domains. They do not establish minimum circuit
complexity except where a separate analytical lower bound is given.
In particular, the tensor checker does not implement or test the
Fomin--H\o ie asymptotic pathwidth theorem; it checks the local
reductions and contraction on supplied decompositions. The $1/5$
rate follows analytically from that cited theorem and the proved
conversion, not from the finite graphs tested.
The tensor checker evaluates Boolean projections. The integer-counting
corollary is justified analytically. The expanded walkthrough adds finite
integer-table checks, while the search and generation validator uses explicit
finite counting oracles. These Python suites are distinct from the later
Rust implementation in \texttt{rs/}, which implements the complete
automatic counting route, including constructive layout. Its
\texttt{ALGORITHM.md} records the fixed-parameter resource argument
and the imported graph lemmas. The structured parity-substitution
experiment in \texttt{EXPERIMENT.md} retains a small kernel while the
input count grows; it illustrates that mechanism and does not measure
the worst-case exponent. Configured resource caps return errors when
the selected route exceeds its limit.

Two further retained suites support the pedagogical expansion. The walkthrough
checks all 4,096 triples of binary gate tables on all eight input assignments
for unique internal extension; all 27 partial pinnings of the running example;
equality splitting at degrees three through six; all 256 pairs of binary
$2\times2$ tables under joint parallel contraction; and exact loop and
frontier fixtures. The query suite checks every Boolean relation on zero
through three bits: 278 relations, 1,844 prefix partitions, and 1,061
satisfying assignments under rank and unrank. It also checks 1,617 invalid
ranks or nonsolutions and exact rejection probabilities for solution counts
one through sixteen. These are finite comparisons against explicit
enumeration and rational arithmetic. They neither implement the general
layout theorem nor establish asymptotic performance experimentally.

The new realization checker addresses a specific possible mismatch:
a graph sketch might fail to be an equality splitting of the actual
circuit's incidence graph. It verifies that correspondence and the
legality of the graph reductions. Its finite graph domain contains
every perfect matching disjoint from a fixed Hamiltonian cycle on
$4,6,8,10,12$ labeled vertices, 42 specified bridged block chains,
and 32 seeded 2-connected cubic graphs. The 3,729 graph checks are
distinct from the seven semantic checks. The general realization and
critical-structure theorems are justified by their analytical proofs.

The user's suggestion of a cheaper aggregate led to the
characteristic-two fingerprint investigation in
Appendix~\ref{sec:aggregation}. It was assessed analytically: a
polynomial zero bound, a simultaneous hardwiring argument, exact
cycle accounting, and a linear-algebra obstruction suffice. Separate
mathematical and source reviews checked this appendix. No additional
experiment was needed, and no improved general exponent resulted.

\paragraph{Reproducibility.}
The manuscript and research notes are distributed under
\href{https://creativecommons.org/licenses/by/4.0/}{CC BY 4.0};
the code and build automation use the MIT license. The
\href{https://github.com/SamuelSchlesinger/understanding-nondeterminism}{public repository}
contains the license texts, citation metadata, and automated checks.

The entry point is \texttt{main.tex}. It includes the files in
\texttt{sections/}. The file \texttt{references.bib} has an individual entry
for each cited work. The \texttt{research/} directory contains the
full statement matrix, supporting derivations, examples, and scripts.
The review findings are retained in \texttt{review-notes/}.
\begin{samepage}
Run
\begin{quote}\ttfamily
make check\\
make pdf
\end{quote}
\end{samepage}
from the project root. The check scripts use only the Python standard
library. The PDF uses the installed \LaTeX{} packages declared in
the source. Seeds and exact finite domains are recorded in the scripts.
Compilation checks references and typography mechanically; rendered
pages are inspected separately because a successful build does not
establish legibility. The research record distinguishes completed
checks from pending work and retains the unresolved literature
comparison described in Section~\ref{sec:prior-work}.
